\documentclass[aspectratio=169,11pt]{beamer}

\usetheme{Madrid}
\usecolortheme{default}
\setbeamertemplate{navigation symbols}{}
\setbeamertemplate{headline}{}
\setbeamertemplate{footline}[frame number]
\setbeamertemplate{itemize items}[circle]

\usepackage[T1]{fontenc}
\usepackage[utf8]{inputenc}
\usepackage{lmodern}
\usepackage{amsmath,amssymb,bm}
\usepackage{xcolor}
\usepackage{array}

\definecolor{matrixblue}{RGB}{35,92,145}
\definecolor{matrixred}{RGB}{165,55,55}
\definecolor{matrixgreen}{RGB}{45,125,80}
\definecolor{matrixorange}{RGB}{190,115,25}

\setbeamercolor{structure}{fg=matrixblue}
\setbeamerfont{frametitle}{size=\large}

\title{Basic Mathematics}
\subtitle{Lecture 1: Matrices}
\author{Damian Chorążkiewicz}
\date{}

\newcommand{\R}{\mathbb{R}}
\newcommand{\blue}[1]{\textcolor{matrixblue}{#1}}
\newcommand{\red}[1]{\textcolor{matrixred}{#1}}
\newcommand{\green}[1]{\textcolor{matrixgreen}{#1}}
\newcommand{\orange}[1]{\textcolor{matrixorange}{#1}}
\newcommand{\keyword}[1]{\textbf{#1}}

\begin{document}

\begin{frame}
  \titlepage
\end{frame}

\begin{frame}{Course position}
  \centering
  \Large
  Matrices $\longrightarrow$ Determinants $\longrightarrow$ Inversion $\longrightarrow$ Linear systems

  \vspace{0.9cm}
  \Large
  Today: \keyword{read} $\cdot$ \keyword{combine} $\cdot$ \keyword{multiply} $\cdot$ \keyword{see structure}
\end{frame}

\begin{frame}{Three lecture blocks}
  \centering
  \renewcommand{\arraystretch}{1.6}
  \begin{tabular}{ccl}
    \textbf{I} & 45 min & Reading matrices and basic operations \\
    \textbf{II} & 45 min & Matrix multiplication \\
    \textbf{III} & 45 min & Special structure and row operations
  \end{tabular}
\end{frame}

% =============================================================================
% PART I
% =============================================================================

\begin{frame}[plain]
  \centering
  \vfill
  {\Huge Part I}

  \vspace{0.5cm}
  {\LARGE Reading matrices}
  \vfill
\end{frame}

\begin{frame}{A matrix is more than a table}
  \centering
  \Large
  \[
    \boxed{\text{data}}
    \qquad
    \boxed{\text{coefficients}}
    \qquad
    \boxed{\text{relations}}
    \qquad
    \boxed{\text{transformation}}
  \]

  \vspace{0.8cm}
  \Huge
  matrix $=$ organised information
\end{frame}

\begin{frame}{General form}
  \Large
  \[
    A=
    \begin{pmatrix}
      a_{11} & a_{12} & \cdots & a_{1n}\\
      a_{21} & a_{22} & \cdots & a_{2n}\\
      \vdots & \vdots & \ddots & \vdots\\
      a_{m1} & a_{m2} & \cdots & a_{mn}
    \end{pmatrix}
    \in \R^{m\times n}.
  \]

  \vspace{0.4cm}
  \centering
  \blue{$m$ rows}\qquad \red{$n$ columns}
\end{frame}

\begin{frame}{Dimensions are part of the object}
  \centering
  \Huge
  \[
    A\in\R^{\blue{m}\times\red{n}}
  \]

  \vspace{0.8cm}
  \Large
  shape first, calculation second
\end{frame}

\begin{frame}{Index discipline}
  \centering
  \Huge
  \[
    a_{\blue{i}\red{j}}
  \]

  \vspace{0.5cm}
  \Large
  \blue{first index: row}\qquad \red{second index: column}
\end{frame}

\begin{frame}{Read before calculating}
  \Large
  \[
    M=
    \begin{pmatrix}
      4 & -2 & 0\\
      7 & 1 & 5
    \end{pmatrix}
  \]

  \vspace{0.5cm}
  \[
    M\in\R^{2\times3},
    \qquad
    m_{11}=4,
    \quad
    m_{12}=-2,
    \quad
    m_{23}=5.
  \]
\end{frame}

\begin{frame}{Rows and columns}
  \Large
  \[
    M=
    \begin{pmatrix}
      \blue{4} & \blue{-2} & \blue{0}\\
      7 & 1 & 5
    \end{pmatrix}
  \]

  \vspace{0.4cm}
  \begin{columns}
    \begin{column}{0.48\textwidth}
      \centering
      first row
      \[
        \begin{pmatrix}4&-2&0\end{pmatrix}
      \]
    \end{column}
    \begin{column}{0.48\textwidth}
      \centering
      third column
      \[
        \begin{pmatrix}0\\5\end{pmatrix}
      \]
    \end{column}
  \end{columns}
\end{frame}

\begin{frame}{Common shapes}
  \begin{columns}[T]
    \begin{column}{0.24\textwidth}
      \centering
      \keyword{row}
      \[
        \begin{pmatrix}a&b&c\end{pmatrix}
      \]
      $1\times3$
    \end{column}
    \begin{column}{0.24\textwidth}
      \centering
      \keyword{column}
      \[
        \begin{pmatrix}a\\b\\c\end{pmatrix}
      \]
      $3\times1$
    \end{column}
    \begin{column}{0.24\textwidth}
      \centering
      \keyword{square}
      \[
        \begin{pmatrix}a&b\\c&d\end{pmatrix}
      \]
      $2\times2$
    \end{column}
    \begin{column}{0.24\textwidth}
      \centering
      \keyword{rectangular}
      \[
        \begin{pmatrix}a&b&c\\d&e&f\end{pmatrix}
      \]
      $2\times3$
    \end{column}
  \end{columns}
\end{frame}

\begin{frame}{One array, different meanings}
  \Large
  \[
    A=\begin{pmatrix}2&1\\-1&3\end{pmatrix}
  \]

  \vspace{0.4cm}
  \begin{columns}[T]
    \begin{column}{0.32\textwidth}
      \centering
      data
    \end{column}
    \begin{column}{0.32\textwidth}
      \centering
      coefficients
    \end{column}
    \begin{column}{0.32\textwidth}
      \centering
      map $\bm{x}\mapsto A\bm{x}$
    \end{column}
  \end{columns}
\end{frame}

\begin{frame}{The first question}
  \centering
  \Huge
  What is the shape?

  \vspace{0.8cm}
  \Large
  Only then: which operation makes sense?
\end{frame}

\begin{frame}{Addition: compatibility}
  \centering
  \Large
  \[
    A+B\quad\text{exists only if}\quad
    \operatorname{size}(A)=\operatorname{size}(B).
  \]

  \vspace{0.7cm}
  \[
    (m\times n)+(m\times n)=m\times n
  \]
\end{frame}

\begin{frame}{Addition: entry by entry}
  \Large
  \[
    A=(a_{ij}),
    \qquad
    B=(b_{ij})
  \]
  \[
    A+B=(a_{ij}+b_{ij})
  \]

  \vspace{0.5cm}
  \centering
  same position $\longleftrightarrow$ same position
\end{frame}

\begin{frame}{Addition example}
  \Large
  \[
    \begin{pmatrix}2&-1\\0&3\end{pmatrix}
    +
    \begin{pmatrix}5&4\\-2&1\end{pmatrix}
    =
    \begin{pmatrix}
      2+5 & -1+4\\
      0+(-2) & 3+1
    \end{pmatrix}
    =
    \begin{pmatrix}7&3\\-2&4\end{pmatrix}.
  \]
\end{frame}

\begin{frame}{Subtraction and scalar multiplication}
  \Large
  \[
    A-B=(a_{ij}-b_{ij})
  \]

  \vspace{0.4cm}
  \[
    \lambda A=(\lambda a_{ij})
  \]

  \vspace{0.5cm}
  \centering
  both preserve the matrix shape
\end{frame}

\begin{frame}{Scalar example}
  \Large
  \[
    A=\begin{pmatrix}2&-1\\0&3\end{pmatrix}
  \]
  \[
    -3A=
    \begin{pmatrix}
      -6&3\\
      0&-9
    \end{pmatrix}.
  \]
\end{frame}

\begin{frame}{Linear combinations}
  \Large
  \[
    \alpha A+\beta B
  \]

  \vspace{0.5cm}
  \[
    2\begin{pmatrix}1&0\\-1&2\end{pmatrix}
    -
    \begin{pmatrix}3&1\\2&0\end{pmatrix}
    =
    \begin{pmatrix}-1&-1\\-4&4\end{pmatrix}.
  \]
\end{frame}

\begin{frame}{Familiar laws survive}
  \Large
  \[
    A+B=B+A
  \]
  \[
    (A+B)+C=A+(B+C)
  \]
  \[
    \lambda(A+B)=\lambda A+\lambda B
  \]

  \vspace{0.4cm}
  \centering
  because these operations are entrywise
\end{frame}

\begin{frame}{Defined or not?}
  \Large
  \begin{columns}
    \begin{column}{0.48\textwidth}
      \centering
      \green{defined}
      \[
        (2\times3)+(2\times3)
      \]
      \[
        5(4\times2)
      \]
    \end{column}
    \begin{column}{0.48\textwidth}
      \centering
      \red{not defined}
      \[
        (2\times3)+(3\times2)
      \]
      \[
        (1\times4)-(4\times1)
      \]
    \end{column}
  \end{columns}
\end{frame}

\begin{frame}{Part I checkpoint}
  \Large
  \begin{itemize}
    \item size of the matrix;
    \item meaning of $a_{ij}$;
    \item rows versus columns;
    \item compatible entrywise operations.
  \end{itemize}

  \vspace{0.6cm}
  \centering
  \Huge
  Read the object first.
\end{frame}

% =============================================================================
% PART II
% =============================================================================

\begin{frame}[plain]
  \centering
  \vfill
  {\Huge Part II}

  \vspace{0.5cm}
  {\LARGE Matrix multiplication}
  \vfill
\end{frame}

\begin{frame}{Multiplication is different}
  \centering
  \Large
  Addition:
  \[
    \text{position} + \text{matching position}
  \]

  \vspace{0.5cm}
  Multiplication:
  \[
    \blue{\text{row}}\cdot\red{\text{column}}
  \]
\end{frame}

\begin{frame}{Dimension rule}
  \centering
  \Huge
  \[
    (\blue{m}\times\orange{n})(\orange{n}\times\red{p})
    =\blue{m}\times\red{p}
  \]

  \vspace{0.5cm}
  \Large
  inner dimensions match; outer dimensions remain
\end{frame}

\begin{frame}{Dimension examples}
  \Large
  \begin{align*}
    (2\times3)(3\times4)&=2\times4,\\
    (1\times5)(5\times1)&=1\times1,\\
    (4\times2)(2\times1)&=4\times1.
  \end{align*}

  \vspace{0.4cm}
  \centering
  \red{$(2\times3)(2\times4)$ is not defined}
\end{frame}

\begin{frame}{The entry formula}
  \Large
  \[
    C=AB,
    \qquad
    C=(c_{ij})
  \]
  \[
    c_{ij}=\sum_{k=1}^{n}a_{ik}b_{kj}
  \]

  \vspace{0.5cm}
  \centering
  row $i$ of $A$ meets column $j$ of $B$
\end{frame}

\begin{frame}{Row meets column}
  \centering
  \Large
  \[
    \blue{\begin{pmatrix}a_{i1}&a_{i2}&\cdots&a_{in}\end{pmatrix}}
    \red{\begin{pmatrix}b_{1j}\\b_{2j}\\\vdots\\b_{nj}\end{pmatrix}}
    =
    \sum_{k=1}^{n}a_{ik}b_{kj}
  \]

  \vspace{0.5cm}
  one row $\times$ one column $=$ one number
\end{frame}

\begin{frame}{Example setup}
  \Large
  \[
    A=\begin{pmatrix}1&3\\-2&4\end{pmatrix},
    \qquad
    B=\begin{pmatrix}2&0\\5&-1\end{pmatrix}
  \]

  \vspace{0.5cm}
  \[
    AB\in\R^{2\times2}
  \]
\end{frame}

\begin{frame}{First row of the product}
  \Large
  \[
    c_{11}=1\cdot2+3\cdot5=17
  \]

  \vspace{0.6cm}
  \[
    c_{12}=1\cdot0+3(-1)=-3
  \]
\end{frame}

\begin{frame}{Second row of the product}
  \Large
  \[
    c_{21}=(-2)\cdot2+4\cdot5=16
  \]

  \vspace{0.6cm}
  \[
    c_{22}=(-2)\cdot0+4(-1)=-4
  \]
\end{frame}

\begin{frame}{The product}
  \Huge
  \[
    AB=
    \begin{pmatrix}
      17&-3\\
      16&-4
    \end{pmatrix}
  \]
\end{frame}

\begin{frame}{Reverse the order}
  \Large
  \[
    BA=
    \begin{pmatrix}2&0\\5&-1\end{pmatrix}
    \begin{pmatrix}1&3\\-2&4\end{pmatrix}
    =
    \begin{pmatrix}2&6\\7&11\end{pmatrix}.
  \]

  \vspace{0.5cm}
  \centering
  \Huge
  $AB\neq BA$
\end{frame}

\begin{frame}{Why order matters}
  \centering
  \Large
  \[
    \bm{x}
    \xrightarrow{\ B\ }
    B\bm{x}
    \xrightarrow{\ A\ }
    A(B\bm{x})
  \]

  \vspace{0.8cm}
  \[
    \bm{x}
    \xrightarrow{\ A\ }
    A\bm{x}
    \xrightarrow{\ B\ }
    B(A\bm{x})
  \]

  \vspace{0.5cm}
  different order $=$ different process
\end{frame}

\begin{frame}{Sometimes only one order exists}
  \Large
  \[
    A\in\R^{2\times3},
    \qquad
    B\in\R^{3\times1}
  \]
  \[
    AB\in\R^{2\times1}
  \]

  \vspace{0.6cm}
  \centering
  \red{$BA$ is not defined}
\end{frame}

\begin{frame}{A rectangular product}
  \Large
  \[
    \begin{pmatrix}
      1&0&2\\
      -1&3&1
    \end{pmatrix}
    \begin{pmatrix}
      2&1\\
      0&-1\\
      4&2
    \end{pmatrix}
    =
    \begin{pmatrix}
      10&5\\
      2&-2
    \end{pmatrix}.
  \]

  \vspace{0.5cm}
  \centering
  $(2\times3)(3\times2)=2\times2$
\end{frame}

\begin{frame}{Matrix times vector}
  \Large
  \[
    A\in\R^{m\times n},
    \qquad
    \bm{x}\in\R^{n\times1}
  \]
  \[
    A\bm{x}\in\R^{m\times1}
  \]

  \vspace{0.6cm}
  \centering
  input vector $\longrightarrow$ output vector
\end{frame}

\begin{frame}{Columns explain $A\bm{x}$}
  \Large
  Let
  \[
    A=\begin{pmatrix}\bm{a}_1&\bm{a}_2&\cdots&\bm{a}_n\end{pmatrix},
    \qquad
    \bm{x}=\begin{pmatrix}x_1\\x_2\\\vdots\\x_n\end{pmatrix}.
  \]
  Then
  \[
    A\bm{x}=x_1\bm{a}_1+x_2\bm{a}_2+\cdots+x_n\bm{a}_n.
  \]
\end{frame}

\begin{frame}{Columns are images of basis vectors}
  \Large
  \[
    A\bm{e}_1=\bm{a}_1,
    \qquad
    A\bm{e}_2=\bm{a}_2,
    \qquad
    \ldots,
    \qquad
    A\bm{e}_n=\bm{a}_n.
  \]

  \vspace{0.7cm}
  \centering
  columns show what the transformation does to basic directions
\end{frame}

\begin{frame}{Associativity and distributivity}
  \Large
  \[
    (AB)C=A(BC)
  \]
  \[
    A(B+C)=AB+AC
  \]
  \[
    (A+B)C=AC+BC
  \]

  \vspace{0.4cm}
  \centering
  grouping may change; order may not
\end{frame}

\begin{frame}{Matrix powers}
  \Large
  \[
    A^2=AA,
    \qquad
    A^3=AAA,
    \qquad
    A^0=I
  \]

  \vspace{0.6cm}
  \centering
  powers require a square matrix
\end{frame}

\begin{frame}{Multiplication checklist}
  \Large
  \begin{enumerate}
    \item read both shapes;
    \item check inner dimensions;
    \item predict the output shape;
    \item compute row by column;
    \item keep the order fixed.
  \end{enumerate}
\end{frame}

\begin{frame}{Part II checkpoint}
  \centering
  \Huge
  \[
    (m\times n)(n\times p)=m\times p
  \]

  \vspace{0.8cm}
  \Large
  Every output entry is one row--column interaction.
\end{frame}

% =============================================================================
% PART III
% =============================================================================

\begin{frame}[plain]
  \centering
  \vfill
  {\Huge Part III}

  \vspace{0.5cm}
  {\LARGE Visible structure and row operations}
  \vfill
\end{frame}

\begin{frame}{Zero matrix}
  \Large
  \[
    0_{m\times n}=
    \begin{pmatrix}
      0&\cdots&0\\
      \vdots&\ddots&\vdots\\
      0&\cdots&0
    \end{pmatrix}
  \]

  \vspace{0.5cm}
  \[
    A+0=A
  \]
\end{frame}

\begin{frame}{Identity matrix}
  \Large
  \[
    I_n=
    \begin{pmatrix}
      1&0&\cdots&0\\
      0&1&\cdots&0\\
      \vdots&\vdots&\ddots&\vdots\\
      0&0&\cdots&1
    \end{pmatrix}
  \]
\end{frame}

\begin{frame}{Correct identity sizes}
  \Large
  For $A\in\R^{m\times n}$:
  \[
    I_mA=A,
    \qquad
    AI_n=A.
  \]

  \vspace{0.6cm}
  \centering
  left identity: $m\times m$\qquad right identity: $n\times n$
\end{frame}

\begin{frame}{Transpose}
  \Large
  \[
    A=(a_{ij})\in\R^{m\times n}
  \]
  \[
    A^{\mathsf T}=(a_{ji})\in\R^{n\times m}
  \]

  \vspace{0.6cm}
  \centering
  rows become columns
\end{frame}

\begin{frame}{Transpose example}
  \Large
  \[
    \begin{pmatrix}
      1&2&3\\
      4&5&6
    \end{pmatrix}^{\mathsf T}
    =
    \begin{pmatrix}
      1&4\\
      2&5\\
      3&6
    \end{pmatrix}.
  \]

  \vspace{0.5cm}
  \centering
  $2\times3\longrightarrow3\times2$
\end{frame}

\begin{frame}{Transpose rules}
  \Large
  \[
    (A^{\mathsf T})^{\mathsf T}=A
  \]
  \[
    (A+B)^{\mathsf T}=A^{\mathsf T}+B^{\mathsf T}
  \]
  \[
    (AB)^{\mathsf T}=B^{\mathsf T}A^{\mathsf T}
  \]

  \vspace{0.4cm}
  \centering
  transpose reverses multiplication order
\end{frame}

\begin{frame}{Symmetric matrices}
  \Large
  \[
    A^{\mathsf T}=A
  \]

  \vspace{0.4cm}
  \[
    \begin{pmatrix}
      2&-1&4\\
      -1&3&0\\
      4&0&5
    \end{pmatrix}
  \]

  \centering
  mirror symmetry across the main diagonal
\end{frame}

\begin{frame}{Diagonal matrices}
  \Large
  \[
    D=\operatorname{diag}(d_1,d_2,\ldots,d_n)
  \]
  \[
    D=
    \begin{pmatrix}
      d_1&0&\cdots&0\\
      0&d_2&\cdots&0\\
      \vdots&\vdots&\ddots&\vdots\\
      0&0&\cdots&d_n
    \end{pmatrix}.
  \]
\end{frame}

\begin{frame}{Powers of a diagonal matrix}
  \Large
  \[
    D=\operatorname{diag}(3,-1,2)
  \]
  \[
    D^2=\operatorname{diag}(9,1,4)
  \]
  \[
    D^3=\operatorname{diag}(27,-1,8)
  \]

  \vspace{0.3cm}
  \centering
  no mixing between coordinates
\end{frame}

\begin{frame}{Diagonal matrices commute}
  \Large
  \[
    D_1=\operatorname{diag}(a_1,\ldots,a_n),
    \qquad
    D_2=\operatorname{diag}(b_1,\ldots,b_n)
  \]
  \[
    D_1D_2=D_2D_1
    =\operatorname{diag}(a_1b_1,\ldots,a_nb_n).
  \]
\end{frame}

\begin{frame}{Visible structure}
  \Large
  Look for:
  \begin{itemize}
    \item many zeros;
    \item equal rows or columns;
    \item proportional rows or columns;
    \item diagonal or triangular patterns;
    \item symmetry;
    \item repeated blocks.
  \end{itemize}
\end{frame}

\begin{frame}{Elementary row operations}
  \Large
  \begin{enumerate}
    \item swap two rows;
    \item multiply a row by a non-zero number;
    \item add a multiple of one row to another.
  \end{enumerate}

  \vspace{0.6cm}
  \centering
  controlled transformations, not random changes
\end{frame}

\begin{frame}{Row-operation notation}
  \Large
  \[
    R_i\leftrightarrow R_j
  \]
  \[
    R_i\leftarrow \lambda R_i,
    \qquad
    \lambda\neq0
  \]
  \[
    R_i\leftarrow R_i+\lambda R_j
  \]
\end{frame}

\begin{frame}{Row-operation example}
  \Large
  \[
    A=
    \begin{pmatrix}
      1&2&-1\\
      2&4&3\\
      0&1&5
    \end{pmatrix}
  \]

  \vspace{0.4cm}
  \[
    R_2\leftarrow R_2-2R_1
  \]
\end{frame}

\begin{frame}{The transformed matrix}
  \Large
  \[
    (2,4,3)-2(1,2,-1)=(0,0,5)
  \]
  \[
    \begin{pmatrix}
      1&2&-1\\
      2&4&3\\
      0&1&5
    \end{pmatrix}
    \longrightarrow
    \begin{pmatrix}
      1&2&-1\\
      0&0&5\\
      0&1&5
    \end{pmatrix}.
  \]
\end{frame}

\begin{frame}{Three layers of a row operation}
  \centering
  \Large
  \[
    \boxed{\text{arithmetic}}
    \longrightarrow
    \boxed{\text{transformation}}
    \longrightarrow
    \boxed{\text{interpretation}}
  \]

  \vspace{0.8cm}
  the new row reveals hidden structure
\end{frame}

\begin{frame}{Column vectors are matrices}
  \Large
  \[
    \bm{u}=\begin{pmatrix}2\\-1\\4\end{pmatrix}
    \in\R^{3\times1}
  \]
  \[
    \bm{u}^{\mathsf T}=\begin{pmatrix}2&-1&4\end{pmatrix}
    \in\R^{1\times3}
  \]
\end{frame}

\begin{frame}{$\bm{u}^{\mathsf T}\bm{u}$ versus $\bm{u}\bm{u}^{\mathsf T}$}
  \Large
  \begin{columns}
    \begin{column}{0.48\textwidth}
      \centering
      \[
        (1\times3)(3\times1)
      \]
      \[
        \bm{u}^{\mathsf T}\bm{u}\in\R^{1\times1}
      \]
    \end{column}
    \begin{column}{0.48\textwidth}
      \centering
      \[
        (3\times1)(1\times3)
      \]
      \[
        \bm{u}\bm{u}^{\mathsf T}\in\R^{3\times3}
      \]
    \end{column}
  \end{columns}
\end{frame}

\begin{frame}{Same vector, two products}
  \Large
  For $\bm{u}=(2,-1,4)^{\mathsf T}$:
  \[
    \bm{u}^{\mathsf T}\bm{u}=2^2+(-1)^2+4^2=21
  \]

  \vspace{0.5cm}
  \[
    \bm{u}\bm{u}^{\mathsf T}
    =
    \begin{pmatrix}
      4&-2&8\\
      -2&1&-4\\
      8&-4&16
    \end{pmatrix}.
  \]
\end{frame}

\begin{frame}{A reliable workflow}
  \Large
  \begin{enumerate}
    \item identify every matrix shape;
    \item decide which operations exist;
    \item predict the result shape;
    \item use the correct local rule;
    \item exploit visible structure;
    \item interpret the result.
  \end{enumerate}
\end{frame}

\begin{frame}{Common errors}
  \Large
  \begin{columns}[T]
    \begin{column}{0.48\textwidth}
      \red{avoid}
      \begin{itemize}
        \item mixing row and column indices;
        \item adding different shapes;
        \item multiplying entrywise;
        \item swapping multiplication order.
      \end{itemize}
    \end{column}
    \begin{column}{0.48\textwidth}
      \green{replace with}
      \begin{itemize}
        \item read $a_{ij}$ carefully;
        \item check dimensions first;
        \item use row--column products;
        \item preserve the order.
      \end{itemize}
    \end{column}
  \end{columns}
\end{frame}

\begin{frame}{Final questions}
  \Large
  \begin{itemize}
    \item What is the shape?
    \item What does each index mean?
    \item Which operations are defined?
    \item What is the output shape?
    \item Does order matter?
    \item What structure can I see?
  \end{itemize}
\end{frame}

\begin{frame}[plain]
  \centering
  \vfill
  {\Huge Read the shape.}

  \vspace{0.5cm}
  {\Huge Then calculate.}
  \vfill
\end{frame}

\end{document}
